\section{Discussion}

This study presents an automatic, reproducible, and hardware-agnostic approach to establishing a CV system for ant counting and detection. The proposed system is designed to be applicable in most laboratory conditions, enabling researchers to automate the task of counting ants in various experimental setups. A series of experiments characterize the system's robustness across multiple imaging conditions, including variations in lighting and background texture. Detecting ants and other minute arthropods in natural scenes epitomizes the "small-object problem" in CV: even on high-resolution photographs, individual animals may occupy only a handful of pixels, leading to missed detections or spurious false positives. This study integrated state-of-the-art deep learning models with SAHI to address this challenge, successfully detecting densely populated ant images with high correlation to human tallies.

\subsection{Ecological and Applied Implications}

Ecologically, the outcomes of this study have several implications for ant research and the broader field of ant behavior and ecology. Automated, minute-by-minute traffic series open research avenues that were previously impractical. Classic food-particle studies \citep{hooper_food_1997}, where ant preferences were inferred from the number of food fragments removed, can now be expanded from a few dishes to dozens of colonies, allowing investigators to chart how size preferences shift across full diel cycles and microclimatic fluctuations. In route-learning experiments \citep{czaczkes_ant_2013, gruter_collective_2015}, thousands of individually detected trajectories per hour make it possible to separate persistent "scouts" from conservative trail followers and to link those behavioral phenotypes to colony hunger or disturbance. Crucially, comparative reviews show that even closely related species diverge sharply in recruitment modes and trail geometry, reflecting differences in colony size, habitat structure, and prey predictability \citep{dejean_foraging_2024}. These emergent strategies underpin general theories of distributed intelligence, yet empirical datasets remain thin because individual foragers are only millimeters long and move rapidly through cluttered backgrounds. High-throughput detection therefore fills a long-standing data gap, providing the fine-scale observations needed to test how recruitment rules evolve across environments.

The same behavioral stream has direct applied value. Because foraging rates are tightly gated by temperature and humidity, continuous traffic curves pinpoint the narrow environmental window when toxic baits will be collected most efficiently—refining the coarse day-versus-night heuristic used in earlier pest-control work \citep{vogt_timing_2005}. Departures from baseline activity or altered trail geometry can also serve as early-warning indicators of sublethal stress from pesticides or heavy-metal pollution, allowing managers to intervene before outright colony collapse. At landscape scale, patrol-rate maps derived from automated detections can be correlated with herbivore damage or seed-removal assays, strengthening biocontrol and seed-dispersal models without destructive sampling. Finally, the framework lays a foundation for broader ecological informatics. Real-time detections streamed from distributed camera traps supply the data needed to assimilate and recalibrate agent-based colony models \citep{li_chaosorder_2014}, transforming static simulations into living predictive tools. Open release of our annotated image sets will seed transfer-learning efforts for other small arthropods—aphids, springtails, parasitoids—while the low compute footprint makes edge-based deployment feasible in remote or battery-powered observatories. By converting image streams into a continuous, standardized behavioral sensor, our approach links individual movement to population dynamics, biogeochemical fluxes, and ultimately, environmental decision-making.

\subsection{Methodological advances: robustness, scalability, and accessibility}

Our study fills three methodological gaps that have hampered the wider adoption of machine-learning pipelines for insect monitoring. First, this study places repeatability as a focused objective. Most auto-detection studies report accuracy from a single, homogeneous imaging setup, such as fixed arenas with threshold tracking \citep{cartas_ayala_behavior_2016,xu_segmentation_2022,sabattini_anttracker_2023}, individually tagged ants \citep{sclocco_integrating_2021}, or uncontrolled field scenes \citep{zhang_aphid_2023,grijalva_detecting_2024}. Because calibration and evaluation are often performed on the same dataset, performance metrics can be inflated and generalizability left untested \citep{chen_common_2025}. Therefore, we benchmarked seven deliberately contrasting laboratory environments, varying glare, background texture, and sensor angle, and showed how each factor modulates system reliability. This multi-condition design, implemented using simple tools such as petri dishes, cotton thread, a camera, and a consumer-grade GPU costing less than \$2,000, provides practical guidance for ecologists working in heterogeneous environments.
Second, this study directly addresses the "small and dense" problem. By combining modern detectors with SAHI, we accurately detect ants that are both very small and densely packed, a scenario that often defeats standard object detection pipelines. Because SAHI is not limited to specific taxa, the same approach can be applied to other densely clustered targets such as aphids \citep{xu_segmentation_2022,zhang_aphid_2023,grijalva_detecting_2024} or even distant bird flocks, extending its usefulness well beyond ant research.
Third, we demonstrate that the approach is truly hardware agnostic. Using a family of models that differ by a factor of twenty in parameter count, we show that a lightweight network with fewer than three million parameters achieves performance within a few percentage points of its much larger counterpart. This finding is crucial for field ecology, where power availability and equipment weight are major constraints: the small model runs efficiently on an inexpensive Raspberry Pi-class edge device, challenging the perception that deep learning requires high-performance clusters or climate-controlled laboratories. Together, these advances make automated insect detection both robust and widely accessible.

\subsection{Study Limitations and Future Work}

While the proposed CV system establishment procedure demonstrates promising results in both concentrated and sparse ant images, several noteworthy limitations emerged when testing on external datasets. To illustrate these limitations, we evaluated the system using images from the publicly available iNaturalist dataset \cite{horn_inaturalist_2018}. We selected six representative images from the dataset to demonstrate the limitations. Of the six test images, four contained Argentine ants (\emph{Linepithema humile}) (Figure \ref{fig:inat} a, b, c, e) and two featured red harvester ants (\emph{Pogonomyrmex barbatus}) (Figure \ref{fig:inat} d, f). Since these species differ from those used in our training dataset, they provide a valuable test case for assessing model limitations.
The RT-DETR-L model, identified as the best-performing model in Study 1, was trained on over 1,000 ant images from our experimental dataset. Consequently, the model struggled to accurately detect ants in these unfamiliar images, particularly when subjects were larger than expected or blended with complex backgrounds. The model performed well on relatively clear images (Figure \ref{fig:inat}a and d), but struggled with images containing small debris such as sand or leaf shadows, producing numerous false positives (Figure \ref{fig:inat}b and e). A particularly illustrative example appears in Figure \ref{fig:inat}c, where the ants are significantly larger than those in our calibration set, with much more prominent morphological features. Here, the model failed to detect the actual ants but incorrectly identified ant legs as separate ants due to their similar size and coloration. Finally, in Figure \ref{fig:inat}f, where ants are barely visible even to human observers due to background camouflage, the model again failed to achieve detection.
These examples underscore the model's limitations when applied to images that differ substantially from the calibration dataset. While not the primary focus of this study, broader application would require consideration of ant species diversity and morphological variation. Enhanced model robustness and generalizability could be achieved by incorporating a wider range of ant images in the calibration set, preferably spanning multiple species, complex backgrounds, and varied distance scales.
Several approaches could address these limitations. First, a more complex model architecture, such as RT-DETR-X with 34 million additional parameters (nearly 80\% more than RT-DETR-L), might improve generalizability. However, this approach would require higher-grade GPUs with sufficient video memory for model calibration, more extensive training datasets to prevent overfitting, and longer processing times compared to smaller models.
Alternatively, rather than increasing model complexity or dataset size, a data augmentation approach could simulate debris and background complexity within the existing calibration set. For instance, treating images as two-dimensional matrices, we could initialize multiple multivariate Gaussian distributions to model the probability of debris and background complexity, then randomly sample from these distributions to generate dark or light pixel clusters that mimic natural debris. While this approach may seem simplistic, it could provide a valuable starting point for improving model performance from a different methodological perspective.

\begin{figure}[H]
    \centering
    \includegraphics[width=\textwidth]{figure_12.jpg}
    \caption{Ant images from the iNaturalist dataset \cite{horn_inaturalist_2018}. The ant species are \emph{Linepithema humile} (a, b, c, e) and \emph{Pogonomyrmex barbatus} (d, f). The blue bounding boxes indicate the detection results from the RT-DETR-L model calibrated in Study 1 and detection made by Slicing Aided Hyper Inference (SAHI) \cite{akyon_slicing_2022}.}
    \label{fig:inat}
\end{figure}

\textit{Alt text}: External dataset images from iNaturalist \cite{horn_inaturalist_2018}.


\subsection{Data Dissemination and Web-Based Application}

The source code, the dataset organized in YOLO format, along with the calibrated YOLO11n, YOLO11m, and RT-DETR-L model weights (.pt), can be accessed for generating and evaluating the CV system and is publicly available on GitHub (\url{https://github.com/Niche-Lab/ant-detective}). To support researchers and practitioners, a web-based Streamlit application called Ant Detective (\url{https://ant-detective.streamlit.app}) has also been developed. This application allows users to upload images and receive a folder containing detection results in YOLO format and visualized images with blue bounding boxes.
